\section{Metatheory and proof status}
This section separates unconditional results, proofs with explicit premises, and declarations made with \texttt{Conjecture}. A Coq conjecture is an assumption available to later proofs. A theorem ending in \texttt{Qed} can therefore still depend on a conjecture; its assumption audit determines the distinction.

\subsection{Established structural results}
\begin{theorem}[Preservation when the type computes]
For every context $\Gamma$ and terms $t,A,A'$,
\[
\chk{\Gamma}{t}{A}\ \land\ A\longmapsto A'
\quad\Longrightarrow\quad\chk{\Gamma}{t}{A'}.
\]
\end{theorem}
\noindent This follows directly from conversion and \textsc{ch-expand}; it does not require term preservation.
\src{progress/TypeRules.v: preservation_type}

Write $\operatorname{whd}(t,h)$ when $t\Downarrow u$ for some term $u$ with the stable head classification $h$ used in \texttt{Progress.v}. These are the formalization's designated rigid heads; this notation does not classify every possible weak-head term.
\begin{theorem}[Conversion preserves stable head classification]
\[
 t\equiv u\ \land\ \operatorname{whd}(t,h_1)\ \land\ \operatorname{whd}(u,h_2)
 \quad\Longrightarrow\quad h_1=h_2.
\]
\end{theorem}
\src{progress/Progress.v: conv_whd}
\begin{theorem}[Canonical positions are distinct under conversion]
\[
 c\equiv d\ \land\ c\Downarrow_{\rm pos}n\ \land\ d\Downarrow_{\rm pos}m
 \quad\Longrightarrow\quad n=m.
\]
\end{theorem}
\src{progress/Progress.v: conv_pos}
\begin{theorem}[Coverage supplies a convertible label]
\[
 a\in_s L\ \land\ \operatorname{covers}(\Psi,L)
 \quad\Longrightarrow\quad\exists d\in\Psi,\ a\equiv d.
\]
\end{theorem}
\src{progress/Progress.v: spine_covered}
\begin{theorem}[A value application has a sort-classified type]
If $\chk{\varnothing}{f\,a}{T}$, $\operatorname{Val}(f\,a)$, $T\equiv U$, and $\operatorname{whd}(U,h)$, then $h=\mathsf{HSort}$.
\end{theorem}
\noindent Value applications here are the applied fixed-point forms designated by the definition of \texttt{value}.
\src{progress/Progress.v: muapp_sort}

\subsection{Progress and the remaining signature interface}
Define the property
\[
\mathsf{Progress}:\!\iff\forall t,A,\quad
\chk{\varnothing}{t}{A}\Longrightarrow
\bigl(\operatorname{Val}(t)\lor\exists t',\ t\longmapsto t'\bigr).
\]
\begin{theorem}[Progress relative to signature canonical forms]
Assuming \texttt{canonical\_forms\_sig}, $\mathsf{Progress}$ holds.
\end{theorem}
\noindent The theorem \texttt{progress\_proved} has the displayed progress type and depends on the global conjecture \texttt{TypeRules.canonical\_forms\_sig}. It does not discharge that conjecture. The separate declaration \texttt{TypeRules.progress} remains a conjecture.
\src{progress/Progress.v: progress_proved}

A smaller, explicit interface also suffices. Let $|t|$ be \texttt{tsize t}, the source's structural size measure, and define
\[
\mathsf{BP}(N):\!\iff\forall t,A,\quad
|t|\le N\ \land\ \chk{\varnothing}{t}{A}
\Longrightarrow\operatorname{Val}(t)\lor\exists t',\ t\longmapsto t'.
\]
\begin{definition}[Signature tag transport, $\mathsf{STT}$]
For all $N,S_0,i_0,S_1,i_1,c,n,xs,X,\Psi$, assume
\begin{align*}
&\mathsf{BP}(N),\qquad |xs|\le N,\\
&\chk{\varnothing}{xs}{\interp{B_{S_0}(i_0,c)}{X}},\\
&\mu^s S_0\,i_0\equiv\mu^s S_1\,i_1,\\
&c\Downarrow_{\rm pos}n,\qquad
  \forall d\in\Psi,\ \exists m,\ d\Downarrow_{\rm pos}m,\\
&c\in_s L_{S_0}(i_0),\qquad\operatorname{covers}(\Psi,L_{S_1}(i_1)).
\end{align*}
Then
\[
 \bigl(\exists d\in\Psi,\ d\Downarrow_{\rm pos}n\bigr)
 \quad\lor\quad\bigl(\exists xs',\ xs\longmapsto xs'\bigr).
\]
\end{definition}
\begin{theorem}[Explicit transport implies progress]
$\mathsf{STT}\Longrightarrow\mathsf{Progress}$.
\end{theorem}
\noindent This implication has no global axiom dependency. Its premise $\mathsf{STT}$ is still unproved for arbitrary mixed conversion; having a closed proof of the implication does not supply its premise.
\src{progress/SignatureProgressReduction.v: signature_tag_transport, progress_from_signature_tag_transport}

Write $\equiv_f$ for \texttt{fconv}, the equivalence/congruence generated by core computation and function eta without Eq-$\varphi$. Core computation here includes the primitive eliminator rules, not just lambda beta reduction.
\begin{theorem}[Tag transport without Eq-$\varphi$]
If
\begin{align*}
&\mu^s S_0\,i_0\equiv_f\mu^s S_1\,i_1,\qquad c\Downarrow_{\rm pos}n,\\
&\forall d\in\Psi,\ \exists m,\ d\Downarrow_{\rm pos}m,\\
&c\in_s L_{S_0}(i_0),\qquad\operatorname{covers}(\Psi,L_{S_1}(i_1)),
\end{align*}
then $\exists d\in\Psi,\ d\Downarrow_{\rm pos}n$.
\end{theorem}
\src{progress/SignatureTagTransport.v: signature_live_tag_transport_luna}

\subsection{Outstanding metatheoretic statements}
All conjectures in this subsection are declarations in \texttt{TypeRules.v}, not unconditional proved theorems in the integrated development.
\begin{conjecture}[Normalization]
\[
 \chk{\varnothing}{t}{A}\quad\Longrightarrow\quad
 \exists v,\ t\Downarrow v\ \land\ \operatorname{Val}(v).
\]
\end{conjecture}
\noindent This is existence of a finite evaluation to a designated value; it is not a statement of strong normalization for all reduction paths.
\src{progress/TypeRules.v: normalization}

\begin{conjecture}[Signature canonical forms]
If $\mathcal F_{\varnothing}(I,E,S)$, $\chk{\varnothing}{i}{I}$, and $\chk{\varnothing}{M}{\mu^s S\,i}$, then there exist $c,xs,\Phi$ such that
\begin{align*}
&M\Downarrow\op{in}(c,xs),\qquad L_S(i)\Downarrow\Phi,\qquad c\in_s\Phi,\\
&\chk{\varnothing}{xs}{\interp{B_S(i,c)}{(\op{Carrier}_E(S))}}.
\end{align*}
\end{conjecture}
\noindent The index-typing and family-formation premises are part of the statement and must be retained.
\src{progress/TypeRules.v: canonical_forms_sig}

\begin{conjecture}[Preservation, or subject reduction]
\[
 \chk{\Gamma}{t}{A}\ \land\ t\longmapsto t'
 \quad\Longrightarrow\quad\chk{\Gamma}{t'}{A}.
\]
\end{conjecture}
\begin{proposition}[Conditional multistep preservation]
Assuming preservation,
\[
 \chk{\Gamma}{t}{A}\ \land\ t\Downarrow t'
 \quad\Longrightarrow\quad\chk{\Gamma}{t'}{A}.
\]
\end{proposition}
\noindent The multistep proposition is proved by induction on evaluation, with a global dependency on \texttt{TypeRules.preservation}.
\src{progress/TypeRules.v: preservation, preservation_eval}

\begin{conjecture}[Consistency]
Let $\bot:=\Pi X:\Set_0.\,X$. Then
\[
 \forall t,\ \neg\chk{\varnothing}{t}{\bot}.
\]
\end{conjecture}
\begin{conjecture}[Empty enumeration]
$\forall t,\ \neg\chk{\varnothing}{t}{\op{EnumT}\,\op{nilE}}$.
\end{conjecture}
\begin{conjecture}[Empty enabled signature]
For every $M,S,i,A$,
\[
 \chk{\varnothing}{M}{\mu^s S\,i}\ \land\ L_S(i)\Downarrow[\,]_A
 \quad\Longrightarrow\quad\mathsf{False}.
\]
\end{conjecture}
\src{progress/TypeRules.v: consistency, consistency_enum, consistency_empty_sig}

\begin{conjecture}[Soundness of dead-description evidence]
For every $I,D,X,t$, if $\mathcal D_{\varnothing}(I,D,X)$ and $\operatorname{Dead}(D)$, then
\[
 \neg\chk{\varnothing}{t}{\interp{D}{X}}.
\]
\end{conjecture}
\noindent This claim is scoped to well-typed descriptions and carriers.
\src{progress/TypeRules.v: against_sound}

\subsection{Proved limits of erasure arguments}
The proof modules use auxiliary erasures. These are proof constructions, not coercions inserted by the typing rules. Let $\varepsilon_\varphi$ denote \texttt{phi\_erase}, and let $\longrightarrow_c^*$ denote \texttt{rtc cstep}, the auxiliary combined reduction used in the conversion proofs.
\begin{theorem}[Erased sort endpoints do not reflect]
\[
 \neg\Bigl(\forall t,j,\quad
 \varepsilon_\varphi(t)\longrightarrow_c^*\Set_j
 \Longrightarrow t\equiv\Set_j\Bigr).
\]
\end{theorem}
\noindent The source supplies a concrete raw term \texttt{bad\_t} whose erasure reaches $\Set_0$, while the original term is not convertible to any sort. This refutes an unrestricted erasure-reflection shortcut; it is not a counterexample to typed progress.
\src{progress/ErasureCounterexample.v: bad_erased_sort_path, bad_not_conv_sort, phi_erase_sort_endpoint_not_reflectable}

\begin{theorem}[Branch erasure cannot recover enabled labels]
There are explicit raw signature families $S_{\rm full},S_{\rm empty}$ and an index $i$ such that
\begin{align*}
&0_E\in_s L_{S_{\rm full}}(i),\qquad
  \operatorname{covers}([\,],L_{S_{\rm empty}}(i)),\\
&\operatorname{join}_c\bigl(\varepsilon_b(\mu^s S_{\rm full}\,i),
                   \varepsilon_b(\mu^s S_{\rm empty}\,i)\bigr),\\
&\neg\exists d\in[\,],\ d\Downarrow_{\rm pos}0.
\end{align*}
Here $\varepsilon_b$ is \texttt{branch\_erase}, and $\operatorname{join}_c$ means that both terms reach a common term under $\longrightarrow_c^*$.
\end{theorem}
\noindent The witness families have the same constant unit description and respectively a singleton and an empty enabled list. Joinability after erasure supplies insufficient information to recover a live label. The theorem does not assert conversion of the original refined types.
\src{progress/SignatureTagTransport.v: BranchErasureLabelGap.branch_erased_recovery_label_gap}
